% discussion.tex (v14, 2026-10-20): no subheadings (Scientific Reports recommended structure); six paragraphs: mechanism, products, previous studies, maps and use in new regions, limitations, outlook.
\section{Discussion}\label{sec:discussion}

The way physical knowledge enters the model determined whether machine learning improved on a regional Stefan model and whether the improvement survived the transfer to a region without labels. In the physics-residual formulation (equation~(\ref{eq:residual})), the Stefan term carries the response to the thaw-season climate of each year in a bounded form, and the learner corrects the departures of the effective coefficient that soil and terrain impose. Four results support this division of roles: the gain required terrain and soil covariates and did not appear with climate covariates alone; it was largest where the Stefan error was largest; the attribution of the residual learner pointed to terrain in the Lena Delta and to soil organic carbon in Canada; and Stefan pseudo-labels pulled the learner towards the Stefan model without adding accuracy. Direct ML, which has to learn the climate response from the labels, reached lower point errors when all labels of a region were available, but its Lena Delta error exceeded that of the Stefan model by \ndrefLeregDzabs{} cm when the region was withheld, and with 25 labels it was less accurate than the Stefan model in Alaska and Canada. The gain of the residual learner is therefore conditional, and the conditions are themselves a result: it is worth adding a learner to the physics where soil and terrain information is available and where the regional Stefan coefficient misrepresents the sites, and it is worth anchoring the learner to the physics where the target region has few labels, which is the common situation in permafrost regions.

The comparison with the published products on the same cells and years explains why their accuracy differs between regions. The two ML products relate site observations pooled over many years to multi-year predictors, and their validation uses random or site-based splits on their own site sets\cite{Liu2024,Wei2026}. Where their training sites are close to our label cells (Alaska, and Canada for Liu et al., whose training data include an earlier version of the same field observations), the products were as accurate as our model, and the comparison there is not independent. Where the products have no nearby training sites (the Lena Delta, median distance \ndistweiLemed{} km), they assigned values between about 40 and 70 cm regardless of the observed ALT (correlation \nrLeWei{} and \nrLeLiu{} against \nrLeRo{}), overestimated the thin active layers of the northern delta by about 20 cm, and were less accurate than a physics-residual model trained without any Lena Delta label. Three properties of our workflow produce this difference: the climate input of each prediction is the thawing degree-day sum of the stated year accumulated to the observation day, so that a thin active layer in a cold year is not averaged with thick ones; the Stefan anchor transfers the climate response to cells and years outside the training set; and the evaluation scores every method and product on the same site-years with a mask for product training sites, which is what makes the regional ranking interpretable. The process-based ESA CCI product, which is not trained on ALT observations, had larger errors than our model with five local labels at the external CALM sites, which places the two approaches on one scale for the first time.

These results extend previous comparisons in three directions. Gautam et al.\cite{Gautam2025} found at 68 Alaska CALM sites of one year that a random forest lost most of its accuracy on held-out sites whereas the Stefan model retained it; we confirm this ordering over site-years of 2010--2024 in three regions and under region hold-out, and we show that the two approaches need not be alternatives, because a learner that corrects the Stefan model retains its transferability. Studies of physics-guided learning in lakes and rivers reported that process knowledge mainly helps under data scarcity and extrapolation\cite{Read2019,Jia2021,Willard2022}; the label-number and region hold-out experiments give this statement a quantitative form for ALT, with the number of labels at which the gain becomes resolved and the error without target-region labels. Finally, the finding that random cross-validation underestimates the spatial-block error of direct ML by \noptimLeDz--\noptimCaDz{} cm, and that the region hold-out error is larger still, follows the literature on spatial validation\cite{Roberts2017,Ploton2020,MeyerPebesma2022}; its consequence for ALT products is that errors from random or site-based splits do not indicate the accuracy of a product in a region without training observations, and that such regions are where products are most needed.

The yearly maps turn these results into a product with stated conditions of use. Each map is driven by the ERA5-Land climate of one year with the TDD accumulated to 1~September, so that it can be compared with observations or products of the same year, and it carries an area of applicability and a prediction interval. The regional mean anomalies of the maps followed the observed anomalies in the Lena Delta ($r$ = \nanobsrLe), where the observation network is stable, but only weakly in Alaska ($r$ = \nanobsrAk), where the set of observed cells changes between campaigns and the observed range of regional anomalies therefore exceeds the predicted one; the linear trends of the maps describe the response of the model to the climate of 2011--2024 and are not observed trends. For a new region, the label-number experiments indicate what to expect rather than a fixed procedure. Without local labels, the Stefan model with the coefficient of other regions, or the physics-residual model trained on other regions, is the safer starting point. Local labels improve the model only when they represent the local range of soils and terrain: 25 Lena Delta labels reduced the error of the physics-residual model from \ntlrmseLeRszero{} to \ntlrmseLeRstwfive{} cm, and five labels reduced the external CALM error from \nftAllRszero{} to \nftAllRsfive{} cm, whereas 25 Canadian labels increased the error of the recalibrated Stefan model from \ntlrmseCaPzzero{} to \ntlrmseCaPntwfive{} cm because the thaw depths of the Canadian sites span a wide range. The spatial-block evaluation within the new region and the area of applicability of the resulting map are the two checks that distinguish these cases, and both are part of the workflow.

The study has limitations that bound these conclusions. Half of the Alaska labels come from 2014, which prevented a resolved estimate of the effect of the year-specific input in that region, and the Lena Delta labels come from one data set in \nnblkLe{} blocks, which widens the confidence intervals. The labels are observation sites rather than independent 1~km cells; several sites of one transect can lie in one cell, and although the spatial-block design removes the resulting dependence from the validation, the site-year counts overstate the number of independent samples at the map scale. The climate input is monthly and the soil covariates are static, so that short-term weather and changes in soil moisture are not represented. Active layers thicker than about 3~m, mainly in mountains and at the southern margin of the permafrost zone, lie outside the training range, and the external tests showed that the residual learner does not transfer to such conditions without local labels. The comparison of maps with observations in one common year rests on few cells in the Lena Delta and Canada.

Daily climate inputs, remotely sensed snow and soil-moisture information, and labels from a fourth region such as West Siberia, where CALM and campaign data both exist, are the next steps. Because every number in this study is tied to a stated site, year and validation design, the same protocol can be applied to new products and regions as they appear, and the year-labelled maps can be extended year by year as the climate data are released.
